\section{Related Work}
\label{sec:related-work}

The closest related work is clock-bound repair for timed systems and
TarTar \cite{koelbl2019clock}. Clock-bound repair encodes timed diagnostic traces into
linear real arithmetic and uses MaxSMT solving to compute small
modifications to guard and invariant bounds \cite{koelbl2019clock}. TarTar extends this
line into a repair analysis tool and introduces an admissibility check
based on functional equivalence. MPTA-Repair builds on this foundation,
but targets a different setting: multiple properties, multiple
diagnostic traces, and full-property validation in timed-automata
models from practical industrial cases.

A second related direction is parameter synthesis for parametric timed
automata \cite{alur1993parametric}. Tools such as IMITATOR synthesize parameter valuations
that satisfy timing constraints or preserve behavior around a reference
valuation \cite{andre2010imitator}. These techniques are powerful for design-space
exploration and robustness analysis \cite{bendik2021robustness}. MPTA-Repair differs in its
use of diagnostic traces and its focus on producing small syntactic
repairs for faulty models rather than synthesizing a full parameter
region.

A third related direction is model repair and automated program repair
using SAT, MaxSAT, SMT, and counterexample-guided refinement \cite{weimer2009automatically}.
These works provide search and optimization ideas such as fault
localization \cite{jose2011bugassist}, candidate blocking \cite{mechtaev2016angelix}, repair minimization
\cite{nguyen2013semfix}, and abstraction refinement \cite{clarke2000cegar}. MPTA-Repair adapts this
style of reasoning to the semantics of timed automata, where candidate
validity must consider both timing properties and untimed behavior
preservation.

Finally, this work is related to industrial verification and
dependability engineering \cite{leveson2011safer}. Industry-oriented verification often
emphasizes full regression, explainability, traceability, and lessons
learned rather than only theoretical completeness \cite{clarke2018handbook}. MPTA-Repair
follows this perspective by treating repair as an engineering workflow:
diagnostic traces generate candidates, but acceptance depends on full
regression evidence and admissibility.
